% !TEX root = ../main.tex
Language reward models (RMs) are the load-bearing evaluator of modern alignment: one preference model is
reused across RLHF, best-of-$N$ selection, and data filtering, so a bias in that single component propagates
and compounds. We ask whether RMs inherit EU-protected demographic and intersectional biases; whether such
bias is \emph{low-complexity} (a near-linear direction that null-space projection removes at near-zero cost
to accuracy) or \emph{high-complexity} (entangled with legitimate signal); and, crucially, whether removal is
\emph{robust}. Extending the mechanistic reward-shaping method of \citet{fein2026onebias} to protected
attributes, and adding a closed-form concept-erasure test (\leace{} \citep{belrose2023leace} followed by a
non-linear probe, the \taco{} signature of entanglement \citep{taco2024}), we run matched-pair contrastive
experiments across three EU AI Act high-risk domains: credit, hiring, and essay grading. Preliminary results
on a 0.6B RM show that demographic biases are large and behaviourally removable from the reward scalar. On the
model-response arm, where the RM scores a decision \emph{about} an applicant rather than the applicant, the
reasoning-correctness concept that dominates those scores proves \emph{non-linearly recoverable after optimal
linear erasure}: null-space projection fixes the reward, not the representation. Extending that erasure test
to the demographic directions themselves is our immediate next step. On essay grading the reliability-harm metric becomes
interpretable, and a novel \emph{standpoint-credibility} effect emerges, in which the RM rewards an identical
argument more when it is voiced from a marginalized standpoint. We are scaling the study to eleven production
reward models from 0.6B to 70B.
